\section{The \ams{} Algorithm}
\label{sec:algorithm}

\subsection{Overview}

\ams{} wraps \ms{} (G.11) with a Robbins-Monro adaptation loop.
The inner chain is identical to \ms{}: at each step, a uniform draw determines
whether to make a coarse or fine move; coarse moves project the county-level
plan back to the tract level via the rebalancer.
The adaptation loop runs every $\Delta$ steps (default $\Delta = 50$),
measuring the observed coarse acceptance rate in the current window and
updating $\alpha$ by the Robbins-Monro rule.

\subsection{Design Choices}

Three design choices distinguish \ams{} from a naive adaptive wrapper:

\paragraph{Choice 1: $\gamma_0/t^{3/4}$ step-size schedule.}
As established in Section~\ref{sec:background}, this schedule satisfies both
Robbins-Monro convergence conditions ($\sum \gamma_t = \infty$ and $\sum
\gamma_t^2 < \infty$).
The $t^{-3/4}$ decay is slower than the $t^{-1}$ alternative, preserving
adaptivity into the mid-run phase, while avoiding the second-condition
violation of the $t^{-1/2}$ schedule.
For $T = 5{,}000$ steps and $\Delta = 50$, there are $t_{\max} = 100$
adaptation rounds; the step sizes range from $\gamma_0 = 0.10$ (first round)
to $\gamma_0/100^{3/4} \approx 0.0032$ (last round).
The final perturbations are small but nonzero, allowing the alpha to track
slow non-stationarity.

\paragraph{Choice 2: Adaptation interval $\Delta = 50$.}
At the default target $\tau = 0.30$, approximately $\alpha \cdot \Delta = 0.30
\times 50 = 15$ steps in each window are coarse-move attempts.
A window of 15 binary outcomes gives a standard error of
$\sqrt{0.30 \times 0.70 / 15} \approx 0.12$ on the acceptance rate estimate.
This is large relative to the step size $\gamma_t$, but the Robbins-Monro
convergence theorem handles the noise explicitly.
A smaller $\Delta$ (e.g., 10) would give noisier estimates and slower
convergence; a larger $\Delta$ (e.g., 200) would reduce noise but adapt too
infrequently to track non-stationarity.
We find $\Delta = 50$ empirically robust across all states tested.

\paragraph{Choice 3: Clip to $[0.05, 0.95]$.}
Without clipping, the Robbins-Monro iterate could reach $\alpha = 0$ (pure
fine chain) if the coarse acceptance rate is persistently below target.
At $\alpha = 0$, no coarse moves occur, so the coarse acceptance rate signal
is undefined.
The lower clip $0.05$ ensures at least one coarse move is expected per
$1/0.05 = 20$ steps, providing a nonzero signal.
The upper clip $0.95$ prevents the chain from becoming a pure coarse chain,
which would fail frequently on the rebalancer and provide a misleading
acceptance rate (all rejections, driving $\alpha$ back down in an oscillatory
fashion).

\subsection{Formal Pseudocode}

Algorithm~\ref{alg:ams} gives the complete \ams{} procedure.

\begin{algorithm}[t]
\caption{\ams{}($\text{fine\_adj}$, $\text{coarse\_adj}$,
  $\text{fine\_to\_coarse}$, $\text{pop}$, $k$, $\delta$,
  $\tau$, $\alpha_0$, $\Delta$, $\gamma_0$, $T$, $s$, $p$)}
\label{alg:ams}
\begin{algorithmic}[1]
\State $\alpha \gets \alpha_0$ \Comment{initial coarse-move probability; default 0.30}
\State $\tau \gets 0.30$ \Comment{target coarse acceptance rate}
\State $\delta_c \gets 3\delta$ \Comment{coarse balance tolerance}
\State $\mathrm{window} \gets [\,]$ \Comment{coarse accept/reject outcomes}
\State $\mathrm{alpha\_trace} \gets [\,]$
\State $\Pi \gets \text{InitialPlan}()$
\State $\mathrm{EC\_record} \gets [\,]$ \Comment{all visited plan ECs}
\For{$\text{step} = 1$ \textbf{to} $T$}
  \State $u \gets \text{step\_seed}(s, \text{step})$
  \Comment{MSC\_STEP\_ prefix (same as MultiScale)}
  \If{$u < \alpha$}
    \State $\Pi' \gets \text{CoarseMove}(\Pi, \text{coarse\_adj}, \text{fine\_to\_coarse},
           \text{pop}, k, \delta_c)$
    \State $\text{ok} \gets \text{Rebalance}(\Pi', \text{fine\_adj}, \text{pop}, k, \delta,
           \text{max\_iter}=200)$
    \State $\mathrm{window}.\text{push}(\text{ok})$
    \If{$\text{ok}$}
      \State $\Pi \gets \Pi'$
    \EndIf
  \Else
    \State $\Pi \gets \text{FineMove}(\Pi, \text{fine\_adj}, \text{pop}, k, \delta)$
    \Comment{standard ReCom step}
  \EndIf
  \State $\mathrm{EC\_record}.\text{push}(\EC(\Pi))$
  \If{$\text{step} \bmod \Delta = 0$}
    \State $t \gets \text{step} / \Delta$
    \Comment{adaptation round index}
    \State $\gamma_t \gets \gamma_0 / t^{3/4}$
    \Comment{implements \texttt{gamma\_0 / t.powf(0.75)}}
    \State $r \gets \text{mean}(\mathrm{window})$
    \Comment{observed coarse acceptance rate}
    \State $\alpha \gets \text{clip}\!\bigl(\alpha + \gamma_t(r - \tau),\; 0.05,\; 0.95\bigr)$
    \State $\mathrm{alpha\_trace}.\text{push}(\alpha)$
    \State $\mathrm{window} \gets [\,]$
    \Comment{clear window}
  \EndIf
\EndFor
\State $\text{sorted} \gets \text{sort}(\mathrm{EC\_record})$
\State \textbf{return} $\text{plan at percentile } p \text{ of sorted EC}$
\end{algorithmic}
\end{algorithm}

\paragraph{Return value.}
Following the \ms{} convention, \ams{} returns the plan at percentile $p$ of
the visited plans sorted by $\EC$ (default $p = 0.0$, returning the
minimum-$\EC$ plan seen).
The full alpha trace and final $\alpha$ are written to the run manifest.

\subsection{Data Model and Compositor Integration}

\ams{} adds three fields to the \ms{} configuration struct:

\begin{verbatim}
pub struct AdaptiveMultiScaleConfig {
    pub total_steps: usize,
    pub target_accept: f64,    // default: 0.30
    pub initial_alpha: f64,    // default: 0.30
    pub adapt_interval: usize, // default: 50
    pub gamma_0: f64,          // default: 0.10
    pub pop_tolerance: f64,
    pub coarse_tol: f64,       // = 3 * pop_tolerance
    pub p: f64,
    pub base_seed: u64,
    pub chain_idx: u32,
}
\end{verbatim}

\noindent
\texttt{AdaptiveMultiScaleChain} wraps \texttt{MultiScaleChain} by reference,
adding \texttt{alpha} (current $\alpha$), \texttt{coarse\_accept\_window}
(Vec$\langle$bool$\rangle$), and \texttt{alpha\_trace} (Vec$\langle$f64$\rangle$).
All projection, rebalance, and hierarchy logic is inherited unchanged from
\texttt{MultiScaleChain}.

\subsection{CLI and YAML}

\paragraph{CLI.}
\begin{verbatim}
redist state --state TX --year 2020 \
  --search multiscale-adaptive \
  --multiscale-steps 5000 \
  --ms-target-accept 0.30 \
  --ms-adapt-interval 50 \
  --percentile 0.0
\end{verbatim}

Advanced parameters \texttt{--ms-gamma0} (default 0.10) and
\texttt{--ms-initial-alpha} (default 0.30) are available as expert flags for
diagnosing slow convergence; they default to sensible values for all tested
states.

\paragraph{YAML.}
\begin{verbatim}
algorithm:
  structure: prime-factor
  weights: county
  search: multiscale-adaptive
  multiscale_steps: 5000
  ms_target_accept: 0.30
  ms_adapt_interval: 50
  percentile: 0.0
\end{verbatim}

\subsection{Audit Trail}

Every run with \texttt{--search multiscale-adaptive} records the following
fields in \texttt{runs/\{label\}/\{year\}/index.json}:

\begin{verbatim}
"search": "multiscale-adaptive",
"total_steps": 5000,
"target_accept": 0.30,
"initial_alpha": 0.30,
"final_alpha": 0.33,
"adapt_interval": 50,
"gamma_0": 0.10,
"alpha_trace": [0.30, 0.31, 0.33, 0.33, ...],
"fine_acceptance_rate": 0.73,
"coarse_acceptance_rate": 0.29,
"selected_step": 2847,
"step_seed_formula": "SHA-256('MSC_STEP_' || step:u64le || '_' ||
    chain_idx:u32le || '_' || base_seed:u64le)"
\end{verbatim}

The $\gamma_0$ field is mandatory (even when it equals the default 0.10)
because \texttt{alpha\_trace} cannot be reproduced without the step-size
schedule.
\texttt{final\_alpha} is the last entry of \texttt{alpha\_trace}.
\texttt{fine\_acceptance\_rate} and \texttt{coarse\_acceptance\_rate} are
computed over all steps, not just the last window.
This ensures \texttt{redist label-verify} can confirm all search parameters
and independently reproduce the alpha trace from the base seed.
